\documentclass{beamer}
\usepackage[utf8]{inputenc}

\usetheme{Madrid}
\usecolortheme{default}
\usepackage{amsmath,amssymb,amsfonts,amsthm}
\usepackage{txfonts}
\usepackage{tkz-euclide}
\usepackage{listings}
\usepackage{adjustbox}
\usepackage{array}
\usepackage{tabularx}
\usepackage{gvv}
\usepackage{lmodern}
\usepackage{circuitikz}
\usepackage{tikz}
\usepackage{graphicx}

\setbeamertemplate{page number in head/foot}[totalframumber]

\usepackage{tcolorbox}
\tcbuselibrary{minted,breakable,xparse,skins}

\definecolor{bg}{gray}{0.95}
\lstset{
    language=C,
    basicstyle=\ttfamily\small,
    keywordstyle=\color{blue},
    stringstyle=\color{orange},
    commentstyle=\color{green!60!black},
    numbers=left,
    numberstyle=\tiny\color{gray},
    breaklines=true,
    showstringspaces=false,
}

%------------------------------------------------------------
\title{5.8.9}
\date{October 3,2025}
\author{EE25BTECH11011 - Navya Banavathu}

\begin{document}

\frame{\titlepage}

\begin{frame}{Question}
 The area of a rectangle gets reduced by $9$ square units, if its length is reduced by $5$ units and breadth is increased by $3$ units. If we increase the length by $3$ units and the breadth by $2$ units, the area increases by $67$ square units. Find the dimensions of the rectangle
\end{frame}

\begin{frame}{Solution}
Let the original length and breadth of the rectangle be $l$ and $b$, respectively.

From the problem:
\begin{align}
(l - 5)(b + 3) &= lb - 9  \\
\Rightarrow 3l - 5b &= 6 
\end{align}
\begin{align}
(l + 3)(b + 2) &= lb + 67 \\
\Rightarrow 2l + 3b &= 61 
\end{align}

Equations (2) and (4) form the system:
\begin{align}
\begin{pmatrix}
3 & -5 \\
2 & 3 \\
\end{pmatrix}
\begin{pmatrix}
l \\
b \\
\end{pmatrix}
=
\begin{pmatrix}
6 \\
61 \\
\end{pmatrix}
\end{align}

\end{frame}



\begin{frame}{Solution}

With the augmented matrix followed by row reduction:

\begin{align}
\left(
\begin{array}{cc|c}
3 & -5 & 6 \\
2 & 3 & 61 \\
\end{array}
\right)
\overset{R_2 \leftarrow R_2 - \frac{2}{3}R_1}{\longrightarrow}
\left(
\begin{array}{cc|c}
3 & -5 & 6 \\
0 & \frac{19}{3} & 57 \\
\end{array}
\right)
\end{align}

\begin{align}
\overset{R_2 \leftarrow \frac{3}{19}R_2}{\longrightarrow}
\left(
\begin{array}{cc|c}
3 & -5 & 6 \\
0 & 1 & \frac{171}{19} \\
\end{array}
\right)
\overset{R_1 \leftarrow R_1 + 5R_2}{\longrightarrow}
\left(
\begin{array}{cc|c}
3 & 0 & \frac{969}{19} \\
0 & 1 & \frac{171}{19} \\
\end{array}
\right)
\end{align}

\begin{align}
\overset{R_1 \leftarrow \frac{1}{3}R_1}{\longrightarrow}
\left(
\begin{array}{cc|c}
1 & 0 & \frac{323}{19} \\
0 & 1 & \frac{171}{19} \\
\end{array}
\right)
\end{align}
\end{frame}

\begin{frame}{Solution}
\begin{align}
\therefore
\begin{pmatrix}
l \\
b \\
\end{pmatrix}
=
\begin{pmatrix}
\frac{323}{19} \\
\frac{171}{19} \\
\end{pmatrix}
=
\begin{pmatrix}
17 \\
9 \\
\end{pmatrix}
\end{align}

The dimensions of the rectangle are:
\begin{center}
\boxed{\text{Length } = 17 \text{ units}, \quad \text{Breadth } = 9 \text{ units}}
\end{center}

\end{frame}


\begin{frame}[fragile]{ Code (C)}
\begin{block}{Part 1}
\begin{lstlisting}[language=C, basicstyle=\ttfamily\small, keywordstyle=\color{blue}, frame=single]
#include <stdio.h>
void solve(int* l, int* b) {
    int i, j;
    for (i = 0; i <= 100; i++) {
        for (j = 0; j <= 100; j++) {
            if ((3 * i - 5 * j == 6) && (2 * i + 3 * j == 61)) {
                printf("Found solution: l = %d, b = %d\n", i, j);
                *l = i;
                *b = j;
                return;
            }
        }
    }
    printf("No solution found.\n");
    *l = -1;
    *b = -1;
}
\end{lstlisting}
\end{block}
\end{frame}


\begin{frame}[fragile]{Python Code}
\begin{block}{Part 1}
\begin{lstlisting}[language=Python, basicstyle=\ttfamily\small, keywordstyle=\color{blue}, frame=single]
import ctypes

# Load DLL (adjust filename/path as needed)
lib = ctypes.CDLL('./cod10.dll')

# Declare function argument and return types
lib.solve.argtypes = [ctypes.POINTER(ctypes.c_int), ctypes.POINTER(ctypes.c_int)]
lib.solve.restype = None

# Prepare ctypes int variables to hold results
l = ctypes.c_int()
b = ctypes.c_int()
\end{lstlisting}
\end{block}
\end{frame}


\begin{frame}[fragile]{Python Code}
\begin{block}{Part 2}
\begin{lstlisting}[language=Python, basicstyle=\ttfamily\small, keywordstyle=\color{blue}, frame=single]
# Call the solve function from the DLL
lib.solve(ctypes.byref(l), ctypes.byref(b))

length = l.value
breadth = b.value

if length == -1 or breadth == -1:
    print("No valid solution found.")
else:
    print(f"Length = {length}, Breadth = {breadth}")
\end{lstlisting}
\end{block}
\end{frame}
\end{document}


